\subsection{Affine dimension, triangular parameters and equidimensional pieces}\label{algebra-proof-025-affine-dimension-triangular-parameters-and-equidimensional-pieces}

Authority: \texttt{algebra.tex} 27648--27919, commit \texttt{a04446e57ec1fbc252a871afcec7752fb2807b14}, SHA-256 \texttt{FA8BB92E58A4F78A2BD01B3B6A4A87DE0A0D279F5DD90641B574DD5FBFFFA4F3}. The entire current interval equals the authority after MC-STK-ERR-1601, together with the captured eighteen-line FAC reference and history at \texttt{proposition-finite-gl-dim-polynomial-ring}. Those additions are retained.

Complete arguments and consequences below are editorial evidence linked to the source, not replacement translations. Minimal wording changes retain the original readings in the ledger. No novelty claim is made.

\subsubsection{The original triangular polynomials and the corrected induction}\label{algebra-proof-025-the-original-triangular-polynomials-and-the-corrected-induction}

Source 27659--27720. Retain \(A=k[x_1,\ldots,x_n]\), its maximal ideal \(\mathfrak m\), the finite extension \(\kappa=A/\mathfrak m\), the original \(\alpha_i\), and the fields \(\kappa_i=k(\alpha_1,\ldots,\alpha_i)\) with \(\kappa_0=k\). Each \(\alpha_i\) is algebraic over \(k\). The algebra generated by finitely many such elements is a finite-dimensional domain over a field and is a field: multiplication by a nonzero element is injective on its finite vector space, hence surjective, giving an inverse. Therefore \(\kappa_i=k[\alpha_1,\ldots,\alpha_i]\), as the source requires.

Let \(P_i(T)\in\kappa_{i-1}[T]\) be the original monic minimal polynomial. Lift each of its coefficients through \(k[x_1,\ldots,x_{i-1}]\to\kappa_{i-1}\), retaining the leading coefficient one and the degree \(d_i=\deg P_i\). This gives the specified \(Q_i(T)\); put \(f_i=Q_i(x_i)\). The evaluation map kills each \(f_i\) and is onto \(\kappa_i\). Suppose the induced \(\psi_i\) is an isomorphism. It identifies \[
 k[x_1,\ldots,x_{i+1}]/(f_1,\ldots,f_{i+1})
 \ \simeq\ \kappa_i[x_{i+1}]/(P_{i+1}(x_{i+1}))
 \ \simeq\ \kappa_{i+1}.
\] The first map sends each original variable to its quotient class and uses \(\psi_i\) on its coefficients; its inverse lifts those coefficient classes. The second map evaluates at \(\alpha_{i+1}\), with kernel exactly the minimal polynomial ideal by Euclidean division. It follows that the map proved injective at the end of the step is \(\psi_{i+1}\). MC-STK-ERR-1601 already makes that correction. Both OCC-12261 and OCC-00542 retain it, without a new operation. The two earlier wording reports concern only the declarations of \(\alpha_i\) and \(\kappa_i\).

The same coefficient maps, leaving the remaining variables fixed, give \[
 A/(f_1,\ldots,f_i)\simeq
 \kappa_i[x_{i+1},\ldots,x_n].
\] Thus each prefix ideal is prime. Each next inclusion is strict, since the monic positive-degree \(P_{i+1}\) is nonzero in that polynomial domain and is not a unit. For \(i=n\) the quotient is exactly \(\kappa\), so the prefix ideal is \(\mathfrak m\). The resulting chain from zero to \(\mathfrak m\) has length \(n\). Localizing keeps all its primes and strict inclusions. The maximal ideal of \(A_{\mathfrak m}\) has \(n\) generators, giving the upper dimension bound \(n\); the chain gives the lower bound. Hence \(A_{\mathfrak m}\) is regular of dimension \(n\). For \(n=0\), \(A=k\), \(\mathfrak m=0\), \(\kappa=k\), and the empty chain has length zero. Every construction above then has its empty interpretation.

\subsubsection{Explicit freeness over the original parameter algebra}\label{algebra-proof-025-explicit-freeness-over-the-original-parameter-algebra}

The triangular construction gives more than the source's dimension statement. Put \(D=\prod_{i=1}^n d_i=[\kappa:k]\), with empty product one. Use new independent variables \(t_i\) solely to prove the actual comparison: \[
 C=k[t_1,\ldots,t_n,X_1,\ldots,X_n]/
       (f_i(X_1,\ldots,X_i)-t_i:1\leq i\leq n).
\] Eliminating the \(t_i\) defines an isomorphism \(C\to A\) by \(X_i\mapsto x_i,\ t_i\mapsto f_i(x_1,\ldots,x_i)\). The inverse sends \(x_i\) to the class of \(X_i\); each composite fixes every named generator using the defining equations.

Regard the same quotient successively as an algebra over \(k[t_1,\ldots,t_n]\). Its \(i\)-th defining equation is monic of degree \(d_i\) in \(X_i\), with coefficients in the preceding algebra. For any ring \(E\), division by a monic degree-\(d\) polynomial in \(E[X]\) gives a unique remainder of degree less than \(d\): subtract its leading monomial successively for existence; a nonzero multiple has the same nonzero leading coefficient as that of its multiplier, so cannot have degree less than \(d\), proving uniqueness. This proof remains valid when \(E\) has zero divisors. Induction consequently makes \(C\) a free \(k[t_1,\ldots,t_n]\)-module with the exact basis \[
 X_1^{b_1}\cdots X_n^{b_n},\qquad 0\leq b_i<d_i.
\] The basis contains \(1\), and therefore the coefficient map is injective: a coefficient annihilating \(1\) has all its basis coordinates zero. Under the isomorphism with \(A\), this proves that the original \(f_i\) are algebraically independent and that \(k[f_1,\ldots,f_n]\subset A\) is finite free of rank \(D\), with basis \(x_1^{b_1}\cdots x_n^{b_n}\). In particular this is an explicit finite faithfully flat map, not a coordinate replacement for \(A\).

For positive integers \(a_1,\ldots,a_n\), base change along \(k[t]\to k[t]/(t_1^{a_1},\ldots,t_n^{a_n})\) gives the actual quotient \(A/(f_1^{a_1},\ldots,f_n^{a_n})\), with \(k\)-basis \[
 f_1^{c_1}\cdots f_n^{c_n}x_1^{b_1}\cdots x_n^{b_n},
 \qquad 0\leq c_i<a_i,\quad 0\leq b_i<d_i.
\] Its \(k\)-dimension is \(D\prod_i a_i\). The radical of the quotient ideal is \(\mathfrak m=(f_1,\ldots,f_n)\): a prime containing all the powers contains all the \(f_i\). Thus the quotient is a finite-dimensional local \(k\)-algebra with residue field \(\kappa\). Its composition factors are all \(\kappa\): each simple module over a commutative local ring is its residue field, and finite \(k\)-dimension ensures a composition series. Additivity of \(k\)-dimension then gives its module length \(\prod_i a_i\).

For each prefix of the powered sequence, the same free base-change description is free over \(k[t]/(t_1^{a_1},\ldots,t_{i-1}^{a_{i-1}})\). Multiplication by \(t_i^{a_i}\) on that coefficient ring is injective, as is seen from its monomial basis in the still unrestricted variable \(t_i\); it is therefore injective on the free module. The final quotient has the nonempty displayed basis. This proves directly that every positive-power sequence \(f_1^{a_1},\ldots,f_n^{a_n}\) is \(A\)-regular and remains so at \(\mathfrak m\). No power or coefficient has been absorbed into a new parameter.

There is also the exact graded comparison \[
 \kappa[T_1,\ldots,T_n]\longrightarrow
 \bigoplus_{r\geq0}\mathfrak m^r/\mathfrak m^{r+1},
 \qquad T_i\longmapsto f_i\bmod\mathfrak m^2.
\] Here a coefficient \(\bar a\in\kappa=A/\mathfrak m\) multiplies a degree-\(r\) class by any lift \(a\in A\); a different lift differs by \(\mathfrak m\), so the class is well defined. Writing \(A=\bigoplus_b k[t]x^b\) by the proved freeness yields \(\mathfrak m^r=\bigoplus_b(t)^r x^b\). For each \(r\), both sides of the comparison have \(k\)-basis indexed by \(b\) and the exponent vectors \(c\) with \(\sum c_i=r\). The displayed map takes one such basis to the other, proving bijectivity in every degree, and it respects multiplication by its definition. No ring section \(\kappa\to A\) is assumed, which matters for inseparable residue extensions.

If \(n\geq1\), the number of degree-\(r\) exponent vectors is \(\binom{n+r-1}{n-1}\). Consequently, for every \(r\geq1\), \[
 \dim_k(A/\mathfrak m^r)=D\binom{n+r-1}{n},
 \qquad
 \operatorname{length}_A(A/\mathfrak m^r)=\binom{n+r-1}{n}.
\] The sum of the degree counts for \(0\leq j<r\) gives the binomial coefficient, by partitioning degree-at-most-\(r-1\) exponent vectors according to their last coordinate (or by the elementary Pascal recurrence). For \(n=0\), \(\mathfrak m=0\), the graded algebra is \(k\) in degree zero and all positive powers give \(A/\mathfrak m^r=k\), of dimension and length one. This is treated directly, without factorials with negative indices.

\subsubsection{Global dimension and the exact prime-interval lengths}\label{algebra-proof-025-global-dimension-and-the-exact-prime-interval-lengths}

Source 27722--27779. Since \(A\) is Noetherian and all its maximal localizations are regular of dimension \(n\), the previously checked finite-global-dimension criterion gives \(\operatorname{gldim}A=n\). Every prime lies below a maximal ideal; localization of the corresponding regular local ring is regular. This gives the claimed regularity at all primes. The field case \(n=0\) has global dimension zero. The eighteen retained FAC lines cite the graded Hilbert-syzygy use and the localized projective-space bound. They remain editorial attribution; the source comparison does not silently alter those passages.

For primes \(\mathfrak p\subset\mathfrak q\) of \(A\), localize at \(\mathfrak q\). This is a regular local domain, hence Cohen--Macaulay. Every saturated chain between zero and its maximal ideal has length \(\dim A_{\mathfrak q}=\operatorname{ht}_A\mathfrak q\). Every saturated chain from zero to \(\mathfrak p\) has length \(\operatorname{ht}_A\mathfrak p\), by localizing at \(\mathfrak p\) and the same theorem. These intervals preserve all primes under localization. Concatenating such a zero-to-\(\mathfrak p\) chain with an arbitrary saturated \(\mathfrak p\)-to-\(\mathfrak q\) chain is saturated. Subtracting the two exact finite lengths gives \(\operatorname{ht}\mathfrak q-\operatorname{ht}\mathfrak p\). For equal endpoints the interval has length zero.

Now let the original finite-type domain be \(S=A/P\). Its primes are the primes of \(A\) containing \(P\), with all intervals preserved. A maximal chain of \(S\) begins at its zero prime and ends at a maximal ideal; the preimage of that maximal ideal is maximal in \(A\), of height \(n\). The interval formula gives every such chain length \(n-\operatorname{ht}_A P\). It follows that \[
 \dim S=\dim S_{\mathfrak m}=n-\operatorname{ht}_A P
\] for every maximal ideal \(\mathfrak m\) of \(S\).

For later use, if \(\mathfrak p\) is any prime of this same domain, with preimage \(P'\subset A\), the identical interval calculation gives \[
 \dim S_{\mathfrak p}=\operatorname{ht}_A P'-\operatorname{ht}_A P,
 \qquad
 \dim(S/\mathfrak p)=n-\operatorname{ht}_A P'.
\] Adding retains the exact identity \(\dim S=\dim S_{\mathfrak p}+\dim(S/\mathfrak p)\). This was proved by the original polynomial presentation; it does not assume that an arbitrary essentially finite-type ring has maximal ideals of the same height.

\subsubsection{Dimension at a point and the closed-specialization minimum}\label{algebra-proof-025-dimension-at-a-point-and-the-closed-specialization-minimum}

Source 27786--27867. Let \(S\) be a finite-type \(k\)-algebra, and let \(Z_i=V(\mathfrak q_i)\) be its finitely many irreducible components. Fix \(x\leftrightarrow\mathfrak p\), put \(I'=\{i:\mathfrak q_i\subset
\mathfrak p\}\), and retain \(T=\bigcup_{i\notin I'}Z_i\), \(U=X\setminus T\). For any open \(W\) containing \(x\) and contained in \(U\), its components are exactly the nonempty \(W_i=Z_i\cap W\), \(i\in I'\). A containment between two of them would, after closure in \(X\), give a containment of the two original components; these are distinct and maximal irreducible subsets.

Each \(W_i\) contains a point closed in \(X\), since the component is the spectrum of a finite-type domain and its closed points are dense. Take such a point \(z\). All generalizations of \(z\) lie in \(W\): if a closed complement contained a generalization, it would also contain its specialization \(z\). Therefore every prime chain in the local ring of \(Z_i\) at \(z\) lies inside \(W_i\), retaining its length \(\dim Z_i\) by the preceding domain result. The reverse inequality follows by taking closures of a strict chain of irreducible closed subsets of \(W_i\) in \(Z_i\); intersecting back with the open set recovers the chain, so strictness is retained. Hence \(\dim W_i=\dim Z_i\), and \[
 \dim_x X=\max_{i\in I'}\dim Z_i.
\] The equality of the dimension of a finite union of components with the maximum follows because the largest member of an irreducible chain is contained in one component; its whole chain lies there.

For a maximal \(\mathfrak m\supset\mathfrak p\), every component through \(x\) also passes through the closed point \(x_0\leftrightarrow\mathfrak m\). The minimal primes of \(S_{\mathfrak m}\) are exactly \(\mathfrak q_i S_{\mathfrak m}\) for the components through \(x_0\). Their quotient local rings are \((S/\mathfrak q_i)_{\mathfrak m/\mathfrak q_i}\), by the exact map \([s/t]\mapsto\bar s/\bar t\). Each has dimension \(\dim Z_i\), so \(\dim S_{\mathfrak m}\) is their maximum and is at least \(\dim_xX\). The nonempty open subset \(V(\mathfrak p)\setminus T\) of \(V(\mathfrak p)\) contains a closed point of \(X\), again by the finite-type Jacobson property. For its maximal ideal \(\mathfrak m\), no extra component occurs. This realizes the minimum and proves all three source equalities. For a closed \(x\), the only maximal ideal above \(\mathfrak p\) is \(\mathfrak p\) itself, proving the subsequent closed-point lemma.

These formulas also give the exact upper-semicontinuity statement: for any integer \(d\), \[
 \{x:\dim_xX\geq d\}=\bigcup_{\dim Z_i\geq d}Z_i,
\] a closed set, and the equality-\(d\) locus is its difference from the corresponding equality with threshold \(d+1\). These statements refer to the source's topological dimension at a point, not the height of a nonmaximal prime.

\subsubsection{The decomposition with its weaker exact hypothesis}\label{algebra-proof-025-the-decomposition-with-its-weaker-exact-hypothesis}

Source 27869--27903. The original proof applies the CM maximal-chain theorem to each \(S_{\mathfrak m}\). Therefore all its minimal-prime components have dimension \(\dim S_{\mathfrak m}\). By the preceding domain result, these are exactly the dimensions of the global components \(Z_i\) containing \(\mathfrak m\). Thus two components meeting at a closed point have equal dimension. If two components meet at all, their nonempty closed intersection contains a closed point (indeed a maximal ideal containing their sum). They therefore have equal dimension. Group the finitely many components by their dimension: \[
 T_d=\bigcup_{\dim Z_i=d}Z_i.
\] Each \(T_d\) is closed, and its complement is a finite union of the other groups, also closed. Distinct groups are disjoint. Every nonempty \(T_d\) has all its components of dimension \(d\).

The exact criterion requires less than Cohen--Macaulayness. For a finite-type \(k\)-algebra \(S\ne0\), the following conditions are equivalent: (1) all local rings at maximal ideals are equidimensional; (2) all local rings at all primes are equidimensional; (3) any two intersecting global irreducible components have equal dimension; (4) the function \(x\mapsto\dim_x\operatorname{Spec}S\) is locally constant; (5) there is the displayed finite decomposition into open-and-closed equidimensional sets \(T_d\).

Here are all implications. Condition (2) implies (1). The argument just given for the source proof uses only local equidimensionality at maximal ideals, so (1) implies (3). Grouping components proves (3) implies (5); the point-dimension formula makes the function equal to \(d\) on \(T_d\), giving (4). If (4) holds and \(z\) belongs to components with generic points \(\eta_i\), then \(\dim_{\eta_i}X=\dim Z_i\). The inverse image of any value of the locally constant function is open and closed. Its closedness forces it to contain every specialization of each point it contains. Since \(z\) specializes every such \(\eta_i\), all these component dimensions must agree. Thus (4) implies (3).

Finally assume (3), and fix a prime \(\mathfrak p\). For each minimal prime \(\mathfrak q_i\subset\mathfrak p\), the domain \(S/\mathfrak q_i\) and the formula in the preceding section give \[
 \dim (S/\mathfrak q_i)_{\mathfrak p/\mathfrak q_i}
   =\dim(S/\mathfrak q_i)-\dim(S/\mathfrak p).
\] All the first terms on the right are equal by (3), since these components meet at \(\mathfrak p\). The second term is independent of \(i\). These quotient local rings are precisely the components of \(S_{\mathfrak p}\), so that local ring is equidimensional. This proves (2), completing the equivalence with all original quotient maps retained.

The source CM hypothesis implies (1), but is stronger. For example retain \(S=k[x,y]/(x^2,xy)\). Its nilradical is \((x)\) and its reduced quotient is the domain \(k[y]\), so its spectrum has just one one-dimensional irreducible component and satisfies (3). At \(\mathfrak m=(x,y)\), \(x\) is nonzero: a polynomial multiplier outside \((x,y)\) has a nonzero constant term, so its product with \(x\) cannot belong to \((x^2,xy)\). Every element of the maximal ideal annihilates \(x\), since \(x^2=xy=0\); hence the local depth is zero, whereas its dimension is one. This local ring is not CM. Thus the weaker criterion has actual examples beyond the source hypothesis.

Finite type must be retained in these equivalences. In Heinrich's original semilocal example verified in group 652, both global components have dimension three, yet the local ring at the original \(M_1\) has components of dimensions two and three. That essentially finite-type example satisfies (3) but fails (1). The earlier evidence retains the full original six-variable presentation; it supplies a concrete boundary, not an authorization to weaken the present finite-type hypothesis.

\subsubsection{Product maps, gaps in the dimensions and empty spectra}\label{algebra-proof-025-product-maps-gaps-in-the-dimensions-and-empty-spectra}

The source invokes the disjoint-decomposition/product comparison. Here it is with the actual nilpotents retained. A clopen set \(T\subset\operatorname{Spec}S\) and its complement are quasi-compact, so write them \(D(I)\) and \(D(J)\) for finitely generated ideals \(I,J\). Disjointness gives \(IJ\subset\sqrt0\), and the covering gives \(I+J=S\). Since \(IJ\) is finitely generated by nilpotents, some power \((IJ)^N\) is zero: choose nilpotence exponents for its finite generators and use the same pigeonhole bound as in the fibre proof. The ideals \(I^N,J^N\) are comaximal, since any prime containing their sum would contain both \(I,J\). Therefore choose actual \(a\in I^N,b\in J^N\) with \(a+b=1\). Their product is zero, so \(a^2=a\), and \(D(a)=T\): a prime outside \(D(I)\) contains \(a\), while on \(D(I)\) it contains \(J\), hence \(b\), so it cannot contain \(a=1-b\). Thus \(e=a\) is the desired idempotent.

For the finitely many nonempty \(T_d\), take these idempotents \(e_d\). For distinct \(d,d'\), their product is idempotent and has empty standard open, so it is nilpotent, hence zero. Their sum is idempotent with standard open all of \(X\); thus its complement is a nilpotent idempotent and is zero. Hence \(\sum e_d=1\). The inverse product maps are \[
 S\longrightarrow\prod_d e_dS,\quad s\longmapsto(e_ds)_d,
 \qquad (s_d)_d\longmapsto\sum_d s_d.
\] Orthogonality and the sum-one identity prove both composites are identities; each factor ring has identity \(e_d\). These are exact maps of the original rings, not just maps of their reductions.

When a dimension between zero and \(\dim S\) does not occur, one may include a zero factor \(S_d=0\). Its spectrum is empty and every assertion about heights of its maximal ideals is vacuous; its dimension remains \(-\infty\), not \(d\). The union statement is interpreted using nonempty \(T_d\), and the source's full product index permits these zero factors. If \(S=0\), there are no components or maximal ideals, the union is empty and the product over the empty set is the zero ring in the category of unital commutative rings. No integer range ending at \(-\infty\) is needed. This is a separate proof clarification, not a proposed alteration of the source's dimension convention.

\subsubsection{Source use and continuation}\label{algebra-proof-025-source-use-and-continuation}

All six reports in 27648--27919 are adjudicated. MC-STK-ERR-1601 is retained. The FAC reference and history are read and preserved exactly; their presence is explicitly captured rather than stripped during comparison. The bounded corpus query has two routing records for the Bhatt--Hochster--Ma title, one TeX and one PDF-primary. Both remain unread for this dependency; there is no PDF fallback and no claim they are distinct mathematical works. Heinrich's already read original TeX supplies the precise finite-type boundary through group 652; no additional reading of that paper is claimed. The triangular free-basis proof, graded maps, component criterion and exact product maps above are established directly from the current source. Continue at authority 27920. Lookahead through 28030 is read but unadjudicated.
